\chapter{Введение}
\label{ch:intro}

Одной из распространённых проблем при производстве микрофлюидных устройств и медицинских аппаратов является нестабильность потока в тонких трубках, которая приводит к потерям давления, снижению эффективности оборудования и даже аварийным ситуациям. 

Решить данную проблему можно, например, путём увеличения диаметра трубок или установки дополнительных компрессоров для повышения давления. Однако такой подход требует значительных материальных затрат, усложняет конструкцию и не всегда применим в условиях ограниченного пространства. Кроме того, он не устраняет саму причину нестабильности потока, а лишь маскирует её, что в долгосрочной перспективе может привести к ещё большим эксплуатационным расходам.

Наиболее эффективным решением является точный расчёт и контроль вязкости воздуха, так как именно этот параметр определяет сопротивление течению газа в тонких трубках. Измерение вязкости позволяет заранее прогнозировать поведение потока, оптимизировать геометрию трубопроводов и выбирать оптимальные режимы работы оборудования. Это не только снижает энергопотребление, но и повышает надёжность систем, обеспечивая стабильность технологических процессов. 

Целью работы являлось измерение вязкости воздуха по течению в тонких трубках.